\begin{center}
\hypertarget{native-066-article048}{}\pdfbookmark[2]{48.}{article-048}48.
\end{center}
Ut igitur corporis coelestis locum in spatio pro quovis temporis momento assignare liceat, sequentia in orbita elliptica nota esse oportebit.

I. Longitudo media pro quodam temporis momento arbitrario, quod \emph{epocha} vocatur: eodem nomine interdum ipsa quoque longitudo designatur. Plerumque pro epocha eligitur initium alicuius anni, scilicet meridies 1. Ianuarii in anno bissextili, sive meridies 31. Decembris anno communi praecedentis.

II. Motus medius inter certum temporis intervallum, e. g. in uno die solari medio, sive in diebus $365$, $365\tfrac{1}{4}$ aut $36525$.

III. Semiaxis maior, qui quidem omitti posset, quoties corporis massa aut nota est aut negligi potest, quum per motum medium iam detur (art. 7): commoditatis tamen gratia uterque semper proferri solet.

IV. Excentricitas.

V. Longitudo perihelii.

\SourceLeaf{067}VI. Longitudo nodi ascendentis.

VII. Inclinatio orbitae.

Haec septem momenta vocantur \emph{elementa} motus corporis.

In parabola et hyperbola tempus transitus per perihelium elementi primi vice fungetur; pro II tradentur quae in his sectionum conicarum generibus motui medio diurno analoga sunt (v. art. 19; in motu hyperbolico quantitas $\lambda k b^{-\frac{3}{2}}$ art. 23). In hyperbola elementa reliqua. perinde retineri poterunt, in parabola vero, ubi axis maior infinitus atque excentricitas $=1$, loco elementi III et IV sola distantia in perihelio proferetur.

\begin{center}
\hypertarget{native-067-article049}{}\pdfbookmark[2]{49.}{article-049}49.
\end{center}
Secundum vulgarem loquendi morem inclinatio orbitae, quam nos a $0$ usque ad $180^0$ numeramus, ad $90^0$ tantum extenditur, atque si angulus orbitae cum arcu $\AscendingNode B$ (Fig. 1) angulum rectum egreditur, angulus orbitae cum arcu $\AscendingNode A$ (qui est illius complementum ad $180^0$) tamquam inclinatio orbitae spectatur; in tali tunc casu addere oportebit, motum esse \emph{retrogradum} (veluti si in figura nostra $E\AscendingNode F$ partem orbitae repraesentat), ut a casu altero ubi motus \emph{directus} dicitur distinguatur. Longitudo in orbita tunc ita numerari solet, ut in $\AscendingNode$ cum longitudine huius puncti in ecliptica conveniat, in directione $\AscendingNode F$ autem \emph{decrescat}; punctum initiale itaque a quo longitudines contra ordinem motus numerantur in directione $\AscendingNode F$ tantundem a $\AscendingNode$ distat, quantum aequinoctium vernale ab eodem $\AscendingNode$ in directione $\AscendingNode A$. Quare in hoc casu longitudo perihelii erit longitudo nodi deminuta distantia perihelii a nodo\raisebox{.4ex}{.} Hoc modo alteruter loquendi usus facile in alterum convertitur, nostrum autem ideo praetulimus, ut distinctione inter motum directum et retrogradum supersedere, et pro utroque semper formulas easdem adhibere possemus, quum usus vulgaris saepenumero praecepta duplicia requirat.
